\chapter{Conception du Système VisionGuard}

\section{Sous-système mécanique}
\subsection{Table d'indexation rotative}
\todo{Décrire la table d'indexation : moteur NEMA17, réduction courroie 3:1, calcul des micropas (240 micropas/veine, 27 micropas/degré à 1/16), précision angulaire.}

\subsection{Mécanisme de positionnement caméra}
\todo{Décrire la rotule ball joint pour la phase PoC, puis le servo DS3218 prévu pour la version complète. Justifier le choix de la démarche progressive.}

\subsection{Structure mécanique}
\todo{Profilés aluminium, supports imprimés 3D, contraintes d'encombrement ($\leq 1$~m²).}

\section{Sous-système de vision}
\subsection{Sélection de la caméra}
\todo{Justifier le choix RPi HQ Camera (IMX296, 1.6 MP, obturateur global, 12-bit ADC) face aux alternatives. Avantage obturateur global pour sujets statiques précis.}

\subsection{Optique}
\todo{Objectif 16 mm C-mount, ouverture f/8--f/11, calcul du champ de vision pour détection $\geq 0.8$~mm.}

\subsection{Éclairage}
\todo{3 LEDs blanches, illumination par le bas des veines. Justifier le choix de l'éclairage bas pour les surfaces internes en céramique matte rose.}

\section{Sous-système électronique et contrôle}
\subsection{Architecture matérielle}
\todo{Décrire les trois nœuds : Laptop (maître), STM32 Nucleo F446RE (contrôle mouvement), RPi5 (acquisition image). Justifier la séparation des rôles.}

\subsection{Driver moteur TMC2208}
\todo{Configuration TMC2208 : Vref = 1.21V (courant RMS $\approx 0.85$A), mode stealthChop, 1/16 microstepping. Justifier le choix par rapport aux alternatives (A4988, DRV8825).}

\subsection{Alimentation}
\todo{Schéma d'alimentation : 220V AC $\rightarrow$ 12V/5A (moteur), 12V $\rightarrow$ 5V (buck converter), 5V/5A USB-C PD (RPi5). Isolation des alimentations.}

\subsection{Assignation des broches STM32}
\begin{table}[H]
\centering
\caption{Assignation des signaux STM32 Nucleo F446RE}
\begin{tabular}{lll}
\toprule
\textbf{Signal} & \textbf{Broche STM32} & \textbf{Connecteur} \\
\midrule
STEP  & PB10 & D6 \\
DIR   & PB5  & D4 \\
UART TX & PA2 & CN10 (USART2) \\
UART RX & PA3 & CN10 (USART2) \\
\bottomrule
\end{tabular}
\end{table}

\section{Sous-système logiciel}
\subsection{Architecture logicielle globale}
\todo{Décrire les trois couches logicielles : firmware STM32 (C/HAL), serveur RPi5 (Python/TCP), orchestrateur Laptop (Python). Schéma des flux de données.}

\placeholderfigure{Flux de données VisionGuard}{Flux de données et séquence de commandes du système VisionGuard}

\subsection{Firmware STM32 — Contrôle mouvement}
\todo{Décrire l'architecture firmware : UART HAL polling, parseur de commandes (HOME/MOV/SPEED), génération STEP/DIR. Justifier HAL polling vs interruptions pour ce cas d'usage.}

\subsection{Serveur RPi5 — Acquisition image}
\todo{Décrire le serveur TCP Python sur RPi5 : écoute commandes CAPTURE/DELETE/QUIT, acquisition picamera2 (1456×1088, XBGR8888), horodatage fichiers, log session JSON.}

\subsection{Orchestrateur Laptop — Coordination système}
\todo{Décrire le script Python maître : envoi commandes UART STM32, envoi commandes TCP RPi5, synchronisation mouvement/capture, interface clavier opérateur.}

\subsection{Protocole de communication UART}
\begin{table}[H]
\centering
\caption{Protocole UART Laptop $\leftrightarrow$ STM32 (115200 baud)}
\begin{tabular}{lll}
\toprule
\textbf{Commande} & \textbf{Description} & \textbf{Réponse} \\
\midrule
\texttt{HOME\textbackslash n}        & Définir position zéro & \texttt{OK\textbackslash n} \\
\texttt{MOV:N\textbackslash n}       & Déplacer N micropas   & \texttt{OK\textbackslash n} \\
\texttt{SPEED:V\textbackslash n}     & Régler vitesse (100--2000 µpas/s) & \texttt{OK\textbackslash n} / \texttt{ERR\textbackslash n} \\
\texttt{READY\textbackslash n}       & Envoyé au démarrage STM32 & — \\
\bottomrule
\end{tabular}
\end{table}

\subsection{Protocole TCP RPi5}
\begin{table}[H]
\centering
\caption{Protocole TCP Laptop $\leftrightarrow$ RPi5}
\begin{tabular}{lll}
\toprule
\textbf{Commande} & \textbf{Description} & \textbf{Réponse} \\
\midrule
\texttt{CAPTURE:good\textbackslash n}   & Capturer image — classe \textit{good}   & \texttt{OK:filename} \\
\texttt{CAPTURE:defect\textbackslash n} & Capturer image — classe \textit{defect} & \texttt{OK:filename} \\
\texttt{DELETE:fname\textbackslash n}   & Supprimer un fichier                    & \texttt{OK} / \texttt{ERR} \\
\texttt{QUIT\textbackslash n}           & Fin de session                          & \texttt{BYE} \\
\bottomrule
\end{tabular}
\end{table}

\section{Module d'intelligence artificielle}
\subsection{Approche de détection}
\todo{Décrire l'approche choisie pour la classification des défauts. Justifier le choix CNN vs approches classiques pour ce type de surface (céramique matte, faible contraste).}

\subsection{Pipeline de collecte de données}
\todo{Décrire la procédure de collecte du dataset : étiquetage \textit{good}/\textit{defect}, commandes clavier, volume cible, visite Pursuit Aerospace.}

\subsection{Architecture du modèle}
\todo{À compléter après entraînement — architecture retenue, hyperparamètres, taille du dataset.}

\section{Conclusion}
\todo{Synthèse de la conception et transition vers le chapitre implémentation/tests.}
